\section{Relation with the Probability of Actual Causality}
\label{sub:pearl_actual_cause_prob}

Pearl's Definition~10.3.5 \citep[Ch.~10]{causalityPearl},
\begin{equation}\label{eq:pearl-paac}
P(\text{caused}(x,y\mid e)) \;=\; \frac{P(U_{xy}\cap U_e)}{P(U_e)},
\qquad U_{xy}=\{\mathbf{u}\colon \mathrm{AC}(x,y;M,\mathbf{u})\},
\end{equation}
is the posterior probability, given the evidence, that we are in a noise state for
which the actual-cause verdict holds, with $\mathrm{AC}(x,y;M,\mathbf{u})$ a chosen
actual-causation predicate --- the causal-beam definition of
\citep[Ch.~10]{causalityPearl} or the Halpern--Pearl AC1--AC3
\citep{actualCausalityHalpern}.\footnote{\label{fn:desert_notebook}The computations
supporting the comparison in this section --- Pearl's original desert-traveller
example and the weak-poison variant --- are collected in the companion notebook
\texttt{docs/source/desert\_traveler.ipynb}, which also includes a framework
cross-check against \texttt{ThinSearchSampler} and a hand-rolled enumeration
matching the framework's output under its current root-only witness pool.}
The same construction underwrites Chockler and Halpern's degree of blame
\citep{chocklerResponsibilityBlameStructuralModel2004}, which substitutes a graded
responsibility share inside the expectation. Because \eqref{eq:pearl-paac} is the
closest construction in the literature to \ourapproach{}'s expectation in
Definition~\ref{def:jointnecsuf}, we now ask in detail how it differs from PCI and
how PCI improves on it.

We proceed in three steps. We first replay Pearl's canonical desert-traveller
illustration and show that PCI, on the original deterministic version, recovers
Pearl's verdicts and rankings. We then exhibit one operational gap ---
PCI's expectation collapses Pearl's existential-over-witnesses step into a single
graded integral --- and one substantive gap, where introducing stochastic mechanism
noise downstream of one of the causes splits the two methods on the \emph{strength}
of attribution. We close with three
further differences that follow from the construction without needing their own
worked example.

\subsection{Pearl's worked example: the desert traveller}

We start with the canonical illustration used by Pearl
\citep[\S 10.3.3]{causalityPearl}.

\begin{example}[Desert traveller]\label{ex:desert_traveler}
A traveller has two enemies. Enemy 2 shoots and empties the canteen
($X{=}1$); enemy 1, unaware, poisons the water in the canteen ($P{=}1$).
The traveller dies ($Y{=}1$). Who is the actual cause of death?
\end{example}

The example involves one source of uncertainty, packed into a binary noise variable
$u$:
\begin{itemize}
    \item $u=0$: the traveller drank the poisoned water before the canteen was
    emptied; the cyanide killed him, and the shooter is irrelevant.
    \item $u=1$: the canteen was empty before he could drink; dehydration killed
    him, and the poisoning is irrelevant.
\end{itemize}

Pearl's beam analysis \citep[Ch.~10]{causalityPearl} (or the Halpern--Pearl AC1--AC3
of \citealp{actualCausalityHalpern}, equivalently) at each deterministic state gives
the actual-cause verdicts:
\begin{itemize}
    \item In $u=1$: $X{=}1$ is the actual cause; $P{=}1$ is not.
    \item In $u=0$: $P{=}1$ is the actual cause; $X{=}1$ is not.
\end{itemize}
So $U_{x,y} = \{u=1\}$ and $U_{p,y} = \{u=0\}$. Say a forensic report says: \emph{no cyanide in the body}, 
 which is incompatible with $u=0$. Then $U_e = \{u=1\}$, and
\[
P(\text{caused}(x,y\mid e)) \;=\;
\frac{P(\{u=1\}\cap\{u=1\})}{P(\{u=1\})} \;=\; 1.
\]
The forensics drives the probability that the shooter was the actual cause to
certainty. A toxicology report finding cyanide would result in a symmetric update for the
poisoner.

\subsection{PCI on the same example}

To put the two methods side by side, we compute the responsibility \ourapproach{}
assigns to the shot on the same model.

The endogenous variables are $X, P, C, D, Y \in \{0,1\}$, where $X$ is the shooter's
action, $P$ the poisoner's action, $C$ the cyanide-killed-him path indicator, $D$ the
dehydration-killed-him path indicator, and $Y$ death. The exogenous binary $u$ encodes
which of the two mechanisms gets the chance to act first. The structural equations,
following \citep[p.~323]{causalityPearl}, are
\[
c = p\,(u' \vee x'), \qquad d = x\,(u \vee p'), \qquad y = c \vee d,
\]
where $u', x', p'$ denote complements.

The factual values are $X=1, P=1, Y=1$, with prior $P(u=0)=P(u=1)=\tfrac{1}{2}$. The
factual values of $(C, D)$ are $(1,0)$ at $u=0$ and $(0,1)$ at $u=1$. The suspect set
is $\mathbf{S} = \{X, P\}$, and we evaluate the singleton candidate cause sets $\{X\}$
and $\{P\}$ one at a time. The witness pool $\mathbf{W} \subseteq \{C, D\}$ gives
four possible active witness sets $\mathbf{T} \in
\{\emptyset, \{C\}, \{D\}, \{C,D\}\}$, weighted uniformly by $\Gamma$. The
non-cause root (i.e.\ $P$ when $\mathbf{C} = \{X\}$, and $X$ when $\mathbf{C} = \{P\}$)
is pinned at its factual value: it is not in the witness pool because the
HP-style specification for this example treats co-roots as part of the context, not
as eligible witnesses. The alternative-value distribution $\Delta$ deterministically
flips $1\to 0$. The AC-aligned kernel
$\Phi(\mathbf{C},\mathbf{T}) = p^\Gamma \cdot \mathbb{I}[Y^s] \cdot Y^n$ reduces, in
this binary deterministic case, to $\Phi > 0$ iff $Y$ flips when $\mathbf{C}$ is set
to its alternative and $\mathbf{T}$ is pinned at its factual value.\footnote{The
current \texttt{ThinSearchSampler} implementation pins only non-deterministic (root)
sites as witnesses, so on the same model it computes the necessity expectation over
a different witness pool ($\{u, X, P\}$) and returns $1/12$ for the singleton cause
rather than $1/8$. The companion notebook checks that this number agrees with a
hand-rolled enumeration of the framework's specification, demonstrating that the
framework is correct for the witness pool it currently supports; closing the
remaining gap to the section's $1/8$ requires extending the framework to pin
mediator sites as witnesses.}

At $u=1$, only $(\{X\}, \{C\})$ flips: the required witness $C=0$ pins the cyanide
path at its dormant value; $X$ is the actual cause. At $u=0$, only $(\{P\}, \{D\})$
flips: the witness $D=0$ pins the dehydration path; $P$ is the actual cause. So
$U_{x,y} = \{u=1\}$ and $U_{p,y} = \{u=0\}$, exactly Pearl's verdicts in his
(10.15)--(10.16).

Under uniform $\Gamma$ over the four $\mathbf{T}$s and the
uniform prior on $u$, conditioning on the ``no cyanide'' forensic report (so $P(u=1\mid e) = 1$),
\[
\mathbb{E}[\Phi\mid\mathbf{C}{=}\{X\}, e] \;=\; \tfrac{1}{4},
\qquad
\mathbb{E}[\Phi\mid\mathbf{C}{=}\{P\}, e] \;=\; 0,
\]
while Pearl's Def.~10.3.5 gives $1$ and $0$. Both methods flip the verdict to
``shooter caused death, poisoner did not''. The sign and ordering match in both
regimes. The magnitudes do not: PCI's $\tfrac{1}{4}$ is the $\Gamma$-fraction of
witnesses passing the necessity test, while Pearl's $1$ is the posterior probability
of a single binary claim. PCI is sensitive to how many context settings there are
and what happens in those contexts, whereas Pearl's metric jumps to one as soon as a
single witness set exists that validates the actual-causality claim. The two readings
answer somewhat different questions, but they ride together on the ranking.

\subsection{PCI expectation vs the AC machinery}

The agreement above is reassuring, but the more interesting observation is what
changes underneath. AC's necessity clause is an \emph{existential}: \emph{there
exists} a witness set $\mathbf{T}$ and an alternative value $c'$ such that the
intervention flips $Y$. To verify the existential, one in principle enumerates the
candidate $(\mathbf{T}, c')$ pairs --- this is the main source of AC's combinatorial
hardness. Pearl's definition then averages the indicator of that already-enumerated
predicate over noise states, inheriting the enumeration unchanged.

\ourapproach{} collapses both moves into one expectation: $\Phi$ is averaged over
the joint $(\mathbf{C}, \mathbf{T}, \mathbf{u})$ under
$\Gamma \otimes P_{\mathbf{U}}$. The existential ``does \emph{some} witness work?''
becomes the quantitative ``\emph{what fraction} of witnesses work?'' --- the
information AC uses only one piece of, exposed as a graded score with sampling-based
estimation.

In the desert traveller this is small change: four candidate witness sets, all
trivially enumerable. The point is the scaling. In a model with $n$ candidate
witness variables, AC's existential demands checking up to $2^n$ subsets per noise
state, and Pearl's posterior over the AC indicator inherits that cost per
evaluation. PCI's expectation can be estimated at a Monte Carlo budget the user
controls, with variance independent of the problem's combinatorial size. The
desert-traveller agreement is the small-$n$ end of a curve whose large-$n$ end is
the computational gap that motivates \ourapproach{} in the first place.

So far we have two conceptual differences that look mainly \emph{operational}.
There are also more substantive cases where the two methods give qualitatively
different answers about who is responsible.

\subsection{Desert traveller with stochastic mechanisms}

The desert traveller is a clean preemption case: at any single noise state, one of
the two enemies is uniquely responsible. PCI's necessity factor tracks this exactly,
as we have just seen. A natural follow-up question is whether PCI's
\emph{sufficiency} factor $Y^s$ adds discrimination beyond necessity once a path is
stochastic. The compact answer for this canonical example is that it does not: a
stochastic version of the desert traveller still leaves $Y^s$ degenerate on the two
correct-cause scenarios. Sufficiency does graded work elsewhere in the paper; here we
explain why this particular example cannot stage that work, so the reader does not
mistake the necessity-only comparison above for a complete picture.

\begin{example}[Desert traveller, weak poison]\label{ex:desert_traveler_weak_poison}
Same enemies as in Example~\ref{ex:desert_traveler}. The cyanide is a small dose
that turns out fatal only with probability $\alpha = 0.1$ (independent exogenous
noise $\xi$); $\alpha = 1$ recovers Pearl's original. The structural equations are
\[
c = p\,(u' \vee x'), \quad v_C = c\cdot \xi, \quad d = x\,(u \vee p'),
\quad y = v_C \vee d.
\]
\end{example}

The new endogenous variable $v_C$ is the ``cyanide-was-fatal'' indicator. In
Pearl's original model the cyanide path firing ($c=1$) was equivalent to the
poison being fatal, so the two collapsed into a single binary $c$. Here we
decouple them: $c$ records whether the poisoned water was drunk, $v_C$ whether
the dose then killed. Concretely, $v_C = c \cdot \xi$, so $v_C = 1$ requires
both that the cyanide path fired and that the dose happened to be fatal. As a
new deterministic mediator, $v_C$ joins $c$ and $d$ in the pool of HP-style
witness candidates, taking the witness pool from $\{c, d\}$ to $\{c, v_C, d\}$.

Under the same observation $X=1, P=1, Y=1$, forensic analysis identifies the
mechanism after the fact:
\begin{itemize}
    \item \emph{Scenario A --- cyanide killed:} $u=0$ (drank early), $\xi=1$ (this
    dose happened to be fatal).
    \item \emph{Scenario B --- dehydration killed:} $u=1$ (drank late), $\xi$
    irrelevant.
\end{itemize}

PCI's necessity factor on the weak-poison variant tracks the forensic verdict in the
expected way (the computations are in the companion notebook): $P$ is the necessary
cause in Scenario A, $X$ in Scenario B. The magnitudes shift slightly relative to
Example~\ref{ex:desert_traveler} because the witness pool now includes the extra
mediator $v_C$ ($2^3 = 8$ subsets instead of $2^2 = 4$), and Scenario B's value for
$X$ sits at $0.4875$ rather than $1/4$ because the noise posterior over $\xi$ adds
configurations in which a larger fraction of witness subsets pass the necessity
test. Pearl's Def.~10.3.5 saturates at $1$ for the correct cause in each scenario.
The two methods produce the same ranking; PCI gives a graded reading, Def.~10.3.5
gives the binary AC indicator averaged over noise.

Sufficiency, however, does not separate the two scenarios on this example. PCI's
sufficiency factor $Y^s$ holds the cause at its factual value, pins the witnesses,
and asks how often the resulting outcome agrees with $Y^\star$. In Scenario A the
factual noise state $(u=0, \xi=1)$ is one in which the cyanide path fires
deterministically; holding $P$ at its factual value with any witness regime keeps
$Y^s = Y^\star = 1$. Symmetrically, in Scenario B the dehydration path fires
deterministically and holding $X$ at factual likewise yields $Y^s = 1$. The
sufficiency factor is degenerate on the correct-cause diagonal because the cause's
own path is the one that already produced the factual outcome; reaffirming the cause
reaffirms what it already did. The discrimination exhibited by Pearl's probability
of sufficiency $\mathrm{PS}$ on this model is a different object: $\mathrm{PS}$
conditions on the cause being \emph{absent} and the outcome being \emph{absent} and
then asks whether intervening to turn the cause on produces $Y = 1$, which gives
$\mathrm{PS}(X) \approx 0.76$ and $\mathrm{PS}(P) = 0.1$ at the population level. PCI
$Y^s$ does not coincide with $\mathrm{PS}$ on this example, and on this example PCI
$Y^s$ does not produce the order-of-magnitude separation $\mathrm{PS}$ shows.

The substantive point about PCI versus Pearl's probability of actual causality
therefore rests on what each construction does \emph{not} have, rather than on a
numerical separation in the desert traveller: PCI's causal impact score has a
distinct sufficiency component that, on examples where it engages, distinguishes
causes that reliably reinstate the outcome from those that do not. Pearl's
Def.~10.3.5 averages a single AC indicator and has no separable sufficiency
component to read off. We refer the reader to the continuous signal-mediation
walkthrough in Section~\ref{sec:signal} and the synthetic archetypes in
Section~\ref{sub:synthetic} for examples that exhibit the sufficiency reading in a
shape the desert traveller cannot.

\subsection{Further differences between Pearl's probability of causation and PCI}


\medskip\noindent\textbf{Breadth of aggregation.}
Def.~10.3.5 aggregates over one dimension: the noise $\mathbf{u}$, conditional on
evidence $e$, for a claim $(X=x, Y=y)$ already fixed. PCI aggregates over four: the
suspect set $\mathbf{C}$ and witness set $\mathbf{T}$ via $\Gamma$, the alternative
value $\mathbf{c}'$ via $\Delta$, and the noise via $P_{\mathbf{U}}$. This is what
supports per-feature comparison --- the same PCI expectation, evaluated over varying
$\mathbf{C}$ inside a fixed $\mathbf{S}$, returns the relative responsibility of each
candidate cause. Def.~10.3.5 cannot perform that comparison without a separate
evaluation per claim, each inheriting the AC-predicate cost discussed above.

\medskip\noindent\textbf{Continuous variables.}
The AC predicates inside Def.~10.3.5 are event-shaped: they read off
$\mathrm{AC}(x,y;M,\mathbf{u})$ for propositional $x, y$, and extensions to
continuous variables require discretising the relevant event before the predicate
applies. PCI's $ci$ function is continuous-native (the absolute-difference impact
score of Section~\ref{sec:causal_impact}), and the synthetic benchmark of
Section~\ref{sub:synthetic} exercises that continuity directly.

\medskip\noindent\textbf{Lineage in Pearl's PN/PS.}
Def.~10.3.5 takes a binary actual-causation predicate --- itself a
structural-equation construction that does not appeal to PN/PS at all --- and lifts
it to a probability through $P(\mathbf{u}\mid e)$; $\mathrm{PN}$, $\mathrm{PS}$, and
$\mathrm{PNS}$ appear nowhere in the construction. PCI, by contrast, generalises
$\mathrm{PN}$ and $\mathrm{PS}$ directly: the necessity factor $Y^n$ is the
counterfactual operation underlying $\mathrm{PN}$, the sufficiency factor $Y^s$ is
the counterfactual operation underlying $\mathrm{PS}$, both indexed by the
cause--witness--alternative--noise tuple, with $\Gamma$ and $\Delta$ exposing the
design choices $\mathrm{PN}/\mathrm{PS}$ makes implicitly. So Pearl's Def.~10.3.5
and \ourapproach{} are not competing implementations of the same idea: they inherit
from different parts of Pearl's own toolkit, with \ourapproach{} playing the role of
a context-sensitive, multi-variable extension of the $\mathrm{PN}/\mathrm{PS}$
machinery rather than a probabilistic wrapper around AC.

\medskip
Pearl's Def.~10.3.5 and \ourapproach{} share the same explanatory target ---
probabilistic responsibility built on top of structural information --- but
instantiate it through different mechanisms. Pearl's definition averages a binary AC
indicator against a posterior over noise; PCI integrates a continuous
necessity--sufficiency kernel against a joint distribution over suspect sets,
witness sets, alternative values, and noise. On clean preemption cases, including
the deterministic and weak-poison desert-traveller variants of this section, the two
methods agree on ranking, and PCI's necessity factor reproduces Pearl's diagnosis at
graded resolution. They diverge in two substantive places that the desert traveller
does not exhibit but other examples in the paper do: (1) the expectation replaces an
existential, with the scaling consequences discussed above, and (2) the sufficiency
factor is separable in PCI but absent from Pearl's definition. The cases where
sufficiency does graded work --- the continuous signal-mediation walkthrough of
Section~\ref{sec:signal}, the synthetic archetypes of
Section~\ref{sub:synthetic}, and the SIR benchmark of
Section~\ref{sec:sir_benchmark} --- are where these constructions answer
recognisably different questions, and the question PCI answers is the one we need
for ML-grade causal attribution.
